\begin{table}[H]
\centering
\caption{Úspešnosť útoku a reakcia detektora podľa scenára a nonce stratégie}
\label{tab:scenario-summary}
\small
\begin{tabular}{l r l r l l}
\toprule
Scenár & \multicolumn{2}{c}{R0/R1} & \multicolumn{3}{c}{R2} \\
\cmidrule(lr){2-3} \cmidrule(lr){4-6}
 & SR & Reakcia & SR & Reakcia & Anomália \\
\midrule
S1: Replay            & 0\%   & \texttt{replay}           & 0\%  & \texttt{quarantined} & \texttt{seq\_reuse} \\
S2: AAD Tamper        & 0\%   & \texttt{decrypt\_failed}  & 0\%  & \texttt{quarantined} & \texttt{nonce\_mismatch} \\
S3a: Nonce Enf.       & 100\% & \texttt{ok}               & 0\%  & \texttt{quarantined} & \texttt{nonce\_mismatch} \\
S3b: Cross-Reader     & 0\%   & \texttt{decrypt\_failed}  & 0\%  & \texttt{quarantined} & \texttt{nonce\_mismatch} \\
S4: Seq Rollback      & 11\%  & \texttt{replay}           & 0\%  & \texttt{quarantined} & \texttt{seq\_rollback} \\
S5: Tag Probe         & 0\%   & \texttt{decrypt\_failed}  & 0\%  & \texttt{quarantined} & \texttt{tag\_fail\_streak} \\
S7: Nonce Tamper      & 0\%   & \texttt{decrypt\_failed}  & 0\%  & \texttt{quarantined} & \texttt{nonce\_mismatch} \\
\bottomrule
\end{tabular}
\\[4pt]
\footnotesize SR = úspešnosť útoku (priemer cez 5~behov). A1 a~A2 vykazujú identické výsledky --- tabuľka agreguje podľa nonce stratégie (R0/R1/R2).
\end{table}
